\scanpage{23}
\mspage{1}
% Paper clip at the upper left.

Soit $X/S$ propre, plat, de présentation finie, $f_*(\underline{\mathcal O}_X)\simeq\underline{\mathcal O}_S$ \unsure{univ.!|univ.} $\underline{\mathrm{Pic}}^0_{X/S}$ un schéma abélien $B$. Soit $A=\underline{\mathrm{Pic}}^0_{B/S}$. On définit alors un \deleted{\unsure{schéma}} $A$-torseur $\underline{\mathrm{Alb}}^1_{X/S}$, et un $S$-morphisme

\[
\psi\colon X\longrightarrow\underline{\mathrm{Alb}}^1_{X/S},
\]

ou ce qui revient au $=$, un homomorphisme

\[
\delta\colon X\times_S X\longrightarrow A\qquad\qquad\delta(x,y)=\psi(y)-\psi(x)
\]

tel que

\[
\delta(x,x)=0\qquad.
\]

Pour ceci, on prend, pour deux sections $x,y$ de $X$ sur $S$

\[
\delta(x,y)=\inserted{\text{classe de }}N_yN_x^{-1}\qquad\text{dans }\underline{\mathrm{Pic}}^0_{B/S}(S)
\]

où $N$ est un module inversible de Weil sur $X\times_S B$ (qui existe grâce à l'existence d'une section de $X$ sur $S$).

On prouve la propriété universelle de $(Q=\underline{\mathrm{Alb}}^1_{X/S},\psi)$ pour les $S$-morphismes de $X$ dans les \deleted{schémas} torseurs sous schémas abéliens (ou plus généralement, sous des schémas en groupes commutatifs quelconques?).

Si on définit $Q'$ par $\delta'=\delta^{-1}$, on trouve un isomorphisme $Q'=Q^{-1}$, $\psi'=\psi^{-1}$, et $(Q',\psi')$ est une autre solution du $=$ Pb universel, canoniquement isomorphe au premier. Donc, lorsqu'on définit $(\underline{\mathrm{Alb}}^1,\psi)$ comme solution du Pb universel, avec $A$ le schéma abélien des translations de $\underline{\mathrm{Alb}}^1$, il y a une \uline{convention de signe} à faire pour choisir un isomorphisme $A\simeq\underline{\mathrm{Pic}}^0_{B/S}$. \uline{Nous prendrons la première convention}
% archive sheet number: 22
